\clearpage
\thispagestyle{empty}
\vspace*{0.85cm}
\hfill {\zihao{5}\songti 学号：\ \underline{\makebox[3.0cm][c]{\StudentID}}}\par
\vspace*{2.15cm}
\begin{center}
  {\zihao{2}\fzxiaobiaosong 东南大学学术型研究生\par}
  \vspace{0.75cm}
  {\zihao{2}\fzxiaobiaosong 学位论文开题报告及论文工作实施计划\par}
\end{center}
\vspace*{3.45cm}
\begin{center}
  \begin{minipage}{11.2cm}
    \coverfield{院（系、所）}{\CollegeName}
    \coverfield{学科（专业）}{\MajorName}
    \coverfield{研究生姓名}{\StudentName}
    \coverfield{学科门类与学位级别}{\DegreeType}
    \coverfield{导师姓名}{\AdvisorName}
    \coverfield{入学年月}{\EntranceDate}
    \coverfield{开题报告日期}{\ProposalDate}
  \end{minipage}
\end{center}
\vfill
\begin{center}
  {\zihao{3} 东\ 南\ 大\ 学\ 研\ 究\ 生\ 院}
\end{center}
\clearpage

\thispagestyle{empty}
\vspace*{3.0cm}
\begin{center}
  {\zihao{2}\fzxiaobiaosong 填\ 表\ 须\ 知}
\end{center}
\vspace{1.2cm}
\begingroup
\zihao{4}
\setstretch{1.65}
\setlength{\parindent}{2em}
1、论文开题报告由研究生本人向审议小组报告并听取意见后，由研究生本人填写此表。\par
2、论文开题报告填写完成后，必须经导师审批，通过后方能提交。\par
3、博士生应在第四学期内、硕士生应在第三学期内完成此开题报告。开题报告经研究生秘书在网上审核确认（硕士生至少半年、博士生至少一年）后方可申请答辩。\par
4、研究生开题前应填写查新报告。查新报告对理、工、医、管等学科博士生作为必要环节。博士生查新工作可委托图书馆负责，也可在完成网络文献检索类研究生课程的学习或参加学校组织的网络文献检索培训后，自行组织查新检索，自行组织查新需要详细文献查新述评作为附件。自行查新报告须经导师审查后由开题报告审核专家组审核签字（或盖章）。硕士生和文科博士生开题查新参考上述办法，不作硬性要求。\par
5、本表一式两份，一份研究生自留放入本人“研究生档案材料袋”；一份由院（系、所）保存并归入院（系、所）研究生教学档案。\par
6、学科门类与学位级别指的是工学（或理学等）博士、硕士。\par
7、本表下载区：http://seugs.seu.edu.cn/3676/list.htm。本表电子文档打印时用 A4 纸张，格式不变，内容较多可以加页。\par
\endgroup
\clearpage

\pagenumbering{arabic}
\setcounter{page}{1}
\thispagestyle{plain}
\noindent{\zihao{3}\heiti 一、学位论文开题报告\par}
\vspace{0.08cm}
\begingroup
\fontsize{10.5pt}{12pt}\selectfont
\setlength{\tabcolsep}{3.2pt}
\renewcommand{\arraystretch}{1.52}
\noindent\makebox[0pt][l]{\tikz[remember picture, overlay]\coordinate (proposalTableTop);}\FrameTableIndent\begin{tabularx}{\FrameTabularWidth}{|>{\centering\arraybackslash}p{2.0cm}|*{9}{>{\centering\arraybackslash}X|}}
  \hline
  \zihao{5}\verticalcell{论 文\\题 目} & \multicolumn{9}{c|}{\ThesisTitle} \tabularnewline \hline
  \zihao{5}\verticalcell{研 究\\方 向} & \multicolumn{9}{c|}{\ResearchDirection} \tabularnewline \hline
  \zihao{5}\verticalcell{题 目\\来 源} & \zihao{5}国家 & \zihao{5}部委 & \zihao{5}省 & \zihao{5}市 & \zihao{5}厂、矿 & \zihao{5}自选 & \zihao{5} 有无合同 & \zihao{5}经费数 & \zihao{5}备注 \tabularnewline \hline
   & & & & & & $\sqrt{}$ & 无 & & \tabularnewline \hline
  \zihao{5}\verticalcell{题 目\\类 型} & \zihao{5}\shortstack[c]{基础\\研究} & \zihao{5}\shortstack[c]{应用\\研究} & \zihao{5}\shortstack[c]{综合\\研究} & \zihao{5}其它 & \multicolumn{5}{c|}{\zihao{5}应用研究} \tabularnewline \hline
  \parbox[c][5.1cm][c]{1.75cm}{\centering\zihao{5}\fontsize{10pt}{11.5pt}\selectfont 论文开题前上网检索情况（硕士生可不作要求；理、工、医、管学科博士生应填写并附查新报告）}
    & \multicolumn{9}{c|}{\zihao{5}无} \tabularnewline \hline
  \multicolumn{10}{|c|}{\bfseries \zihao{5}开题报告内容} \tabularnewline
  \multicolumn{10}{|c|}{\zihao{5}（包括：立题依据及价值、研究内容及方法、可行性分析、预期的成果、预计的困难及解决办法）} \tabularnewline
\end{tabularx}
\drawProposalFirstPageContentFrame
\endgroup
\par\vspace{0.12cm}
\leftskip=0pt
\rightskip=0pt
